%%%%%%%%%%%%%%%%%%%%%%%%%%%%%%%%%%%%%%%%%%%%%%%%%%%%%%%
% SECTION 4: RESULTS 2.1 - DATA SOURCES AND STUDY DESIGN
%%%%%%%%%%%%%%%%%%%%%%%%%%%%%%%%%%%%%%%%%%%%%%%%%%%%%%%
% 
% INPUTS NEEDED: 
%   - Dataset sizes and splits
%   - Dictionary sizes
%   - Standards compliance info
%
% FIGURES: None
% TABLES: None
% DEPENDENCIES: Methods section (detailed implementation)
%
%%%%%%%%%%%%%%%%%%%%%%%%%%%%%%%%%%%%%%%%%%%%%%%%%%%%%%%

\section{Results}\label{sec2}

\subsection{Data sources and study design}\label{subsec1}

Our study utilized three primary data sources for development and validation of the OncoCITE framework. First, we accessed the CIViC database (Clinical Interpretation of Variants in Cancer), which provided 1,300 oncology publications with expert-curated variant-evidence relationships serving as ground truth for training our automated extraction system. Second, we generated a synthetic dataset of 21,000 natural language question--SQL query pairs through systematic template expansion with real CIViC entities, distributed across four complexity tiers to ensure comprehensive coverage of clinical query patterns. Third, we compiled comprehensive terminology dictionaries mapping 14,847 gene aliases from HGNC and Ensembl, 8,923 disease terms from ICD-10 and SNOMED-CT, and 12,156 therapy names including generic, brand, and mechanism class nomenclature. All data processing adhered to HGVS nomenclature standards and maintained compatibility with HL7 FHIR-genomics specifications.

%%%%%%%%%%%%%%%%%%%%%%%%%%%%%%%%%%%%%%%%%%%%%%%%%%%%%%%
% KEY DATA IN THIS SECTION:
%
% THREE PRIMARY DATA SOURCES:
%
% 1. CIViC Database:
%    □ 1,300 oncology publications
%    □ Expert-curated variant-evidence relationships
%    □ Ground truth for extraction training
%    □ 70/30 train/test split (mentioned in 2.2)
%
% 2. Synthetic Q-SQL Dataset:
%    □ 21,000 question-SQL pairs
%    □ Template expansion with real CIViC entities
%    □ Four complexity tiers:
%      - Simple (single-table)
%      - Medium (multi-join)
%      - Complex (aggregations)
%      - Super-Complex (nested subqueries)
%
% 3. Terminology Dictionaries:
%    □ 14,847 gene aliases (HGNC, Ensembl)
%    □ 8,923 disease terms (ICD-10, SNOMED-CT)
%    □ 12,156 therapy names (generic, brand, mechanism class)
%
% STANDARDS COMPLIANCE:
%    □ HGVS nomenclature
%    □ HL7 FHIR-genomics compatibility
%
%%%%%%%%%%%%%%%%%%%%%%%%%%%%%%%%%%%%%%%%%%%%%%%%%%%%%%%

%%%%%%%%%%%%%%%%%%%%%%%%%%%%%%%%%%%%%%%%%%%%%%%%%%%%%%%
% CROSS-REFERENCES:
%
% These numbers must match in:
%   → Abstract (1,300 extractions, 21,000 pairs)
%   → Results 2.2 (1,300 publications)
%   → Results 2.4 (21,000 pairs)
%   → Supplementary Methods
%
%%%%%%%%%%%%%%%%%%%%%%%%%%%%%%%%%%%%%%%%%%%%%%%%%%%%%%%
